分子动力学\textrm{(Molecular Dynamics, MD)}是根据\textrm{Newton}力学，按照体系中每个粒子(原子或分子)受力情况，描述粒子运动随时间变化的模拟计算过程。

\subsection{\ch{Si}的熔融}
\subsubsection{金刚石结构的\ch{Si}}
采用第一原理分子动力学\textrm{(ab-initio MD)}模拟64个原子的金刚石\ch{Si}在\textrm{Nos\'e-Hoover}等温的正则系综\textrm{(canonical ensemble)}\footnote{\textrm{canonical}，汉译作``正则''，出自《楚辞\textperiodcentered 离骚》``皇揽揆余於初度兮，肇锡余以嘉名;~名余曰正则兮，字余曰灵均''，《楚辞章句》注%\upcite{Chucizhangju}
:~``正，平也;~则，法也;~灵，神也;~均，调也。言正平可法则者，莫过于天;~养物均调者，莫神于地。高平曰原，故父伯庸名我为平以法天，字我为原以法地。言己上之能安君，下之能养民也。''意思是说(父辈给屈原赐名、字)``正则''、``灵均''隐喻着某种意义，即平、正是天道的象征，原、均是地的象征。因此``正则''的含义是``符合天道''，与\textrm{canonical}的意思\textrm{of, relating to, or forming a canon}意义一致。}中90\textrm{fs}的动力学行为。

在\textrm{MD}模拟中，指定温度下，原子或分子的运动规律是由经典力学描述的，换言之，每个时间步长里，原子或分子是作为经典粒子处理，其受力与运动服从\textrm{Newdon}第二定律。如果粒子的受力是由第一原理计算得到的，就被称为第一原理\textrm{MD}。考虑到计算考虑的正则系综(\textrm{NVT}系综)，计算过程中保持粒子数\textrm{(N)}不变，体积\textrm{(V)}不变并且温度\textrm{(T)}不变。

计算中的金刚石结构的\ch{Si}的格式为\textrm{.cif}格式\footnote{\textcolor{red}{\textrm{cif}格式}}。\textrm{Python}的开源材料分析模块\textrm{(Python Materials Genomics, Pymatgen)}模块，可以把\textrm{.cif}格式文件转换为\textrm{POSCAR}格式。由于\textrm{MD}模拟的体系较大，很多模拟对象是在惯用晶胞基础上扩展成的超晶胞上完成的，为的是将影响体系动力学主要因素都考虑在内。因此，在实际计算中，仔细估计物理量随体系扩张的收敛条件，是选择晶格大小的重要因素。

使用\textrm{Pymatgen}完成\textrm{cif}格式转换成\textrm{POSCAR}的代码如下：

使用超晶胞完成\textrm{MD}模拟，因为超晶胞的大小决定了晶格振动模式的最大波长，所以分子动力学模拟中的超晶胞必须足够大，这样才能包含全部振动模式，因为\textrm{MD}模拟中，很多物理量是和振动模式有关的。这类计算往往可以通过声子计算估计，或者通过几个不同大小的超晶胞模拟估计。

另一方面，对于计算诸如表面问题中的``吸附-基底''或气体和液体的模拟问题中，模拟的晶胞必须足够大，因为模拟引入了周期边界条件，如果最小重复单元太小，会在周期边界上非物理的原子、分子间相互作用。有必要指出的是，这同样适用于此类体系的弛豫问题。

总之，\textrm{MD}模拟中，所选择的超晶胞必须足够大，为的是保证各态遍历假设\textrm{(ergodic simulation)}，确保体系的全部振动模式都包括在内。

64个原子的模拟对象\textrm{POSCAR}

计算控制参数文件\textrm{INCAR}

$\vec k$空间网格剖分文件\textrm{KPOINTS}

在\textrm{INCAR}文件中，计算控制参数\texttt{ISMEAR}，\texttt{SIGMA}，\texttt{ALGO}，\texttt{PREC}和\texttt{ISYM}是求解\textrm{Kohn-Sham}方程波函数相关的。其余都是和\textrm{MD}模拟有关的而控制参数，\texttt{IBRION}=0即开启\textrm{MD}计算，而\texttt{NSW}和\texttt{POTIM}分别是指定\textrm{MD}模拟的离子步的数目和步长，步长的时间单位是\textrm{fs}。

参数\texttt{MDALGO}设置的是等温条件，这里使用的是\textrm{Nos\'e-Hoover}等温模拟，该等温方式中热浴(环境)和体系的能量交换是通过热惯性\textrm{(thermal inertia)}，即\textrm{Nos\'e}质量完成的。\textrm{Nos\'e}质量参数为\texttt{SMASS}。为完成\textrm{NVT}系综模拟，起始温度(\texttt{TEBEG})和终了温度(\texttt{TEEND})保持相同。而通过参数\texttt{ISIF}=2，确保模拟过程中体积不变。

因为计算采用超晶胞，因此\textrm{KPOINTS}文件中只有$\Gamma$点。

为了节省计算开销，采用\textrm{vasp\_gam}，而且可以算得快些。类似的，标准输出(或者\textrm{OSZICAR})中有关参数的说明：

\begin{itemize}
	\item \texttt{T}:~当前离子步的瞬时温度。
	\item \texttt{E}:~总能，对应\textrm{OUTCAR}中的\texttt{ETOTAL}:~(1)源于离子自由度的势能\texttt{F}，(2)\textrm{Nos\'e-Hoover}等温模拟贡献的势能\texttt{SP}和动能\texttt{SK}，(3)离子的动能\texttt{EK}。
	\item \texttt{F}:~\textrm{DFT}计算的体系自由能(因为展宽因子\texttt{SIGMA}，引入了虚拟电子温度)，对应\textrm{OUTCAR}中的\textrm{TOTEN}:~从\textrm{MD}的角度看，这就是离子自由度的势能部分。
	\item \texttt{E0}:~\textrm{DFT}计算的总能，源于\textrm{DFT}总自由能\texttt{F}扣除熵的贡献，对应\texttt{SIGMA}的贡献为零。
	\item \texttt{EK}:~离子的动能，对应\textrm{OUTCAR}中的\texttt{EKIN}。
	\item \texttt{SP}:~\textrm{Nos\'e-Hoover}等温贡献的势能:~对应\textrm{OUTCAR}中的\texttt{ES}。
	\item \texttt{SK}:~\textrm{Nos\'e-Hoover}等温贡献的动能:~对应\textrm{OUTCAR}中的\texttt{EPS}。
\end{itemize}

值得注意的是，每个离子步温度并不是恒定的常数。因为每个离子步的瞬时温度并非可观测的系综平均值。因为模拟中的等温模拟的要求，并不是每一个时间步长里的瞬时温度都要恒定，而只要求系综平均的温度是固定值即可。在\textrm{OUTCAR}文件中，有提示字符串\textcolor{red}{\textrm{mean temperature}}。如果有足够多的原子，在热力学极限下，这种瞬时温度的涨落就可以忽略。

采用\textrm{py4vasp}作图对能量和时间步长。

根据模拟，可以看到继续\textrm{90fs}和\textrm{2000K}的模拟，金刚石结构的\ch{Si}将会熔融，因此可以得到第一原理计算的原子受力，由此可以得到精确的经典的粒子径迹\textrm{(trajectories)}。

\subsubsection{\ch{Si}的熔融}
对64个原子的金刚石构型\ch{Si}，应用\textrm{2000K}的正则系综，完成第一原理\textrm{MD}模拟，时长分别为\textrm{90fs}，\textrm{180fs}和\textrm{195fs}，得到熔融\ch{Si}的对相关函数。

液态的结构可以通过径向分布函数\textrm{(radial distribution function)}表示，在这里也就是对相关函数。大致地说，对相关函数描述的是以一个粒子为中心，在给定距离范围内找到另一个粒子的概率。更准确地说，是以给定的任意参照粒子为中心，在距离$r$的范围内，局域粒子密度与系统平均密度的比例。可以写作
\begin{displaymath}
	g(r)=\dfrac{\Omega_r}{N^2}\langle\sum_i\sum_{j\neq i}\delta(\vec r-\mathbf{R}_{ij})\rangle
\end{displaymath}
这里$\Omega_r$是实空间体积，$N$是粒子数，$\mathbf{R}_{ij}=\mathbf{R}_i-\mathbf{R}_j$，而$\langle\cdot\rangle$表示取系综平均值。
\begin{displaymath}
	A_{\mathrm{obs}}=\langle A(p(t),q(t))\rangle_{\mathrm{time}}=\lim_{t_{\mathrm{obs}}\rightarrow\infinity}\dfrac1{t_{\mathrm{obs}}}\int_0^{t_{\mathrm{obs}}}A(p(t),q(t))\mathrm{d}t
\end{displaymath}
这里$A_{\mathrm{obs}}$表示可观测的宏观物理量，是对物理量$A(p(t),q(t))$对时间的平均，也就是在$p(t)$和$q(t)$定义的相空间内的点随时间$t$变化的曲线。

晶体中的长程序，在对相关函数形式中，表现为一系列特定的峰值位置。而液态中，粒子与固态中类似，依然受很强的束缚，但体系在发生熔融的时候，对相关函数则有了很大的改变。

应用对相关函数，可以计算任何与原子对有关的函数\textrm{B}。
\begin{displaymath}
	\mathrm{B}_{\mathrm{obs}}=\langle\sum_i\sum_{j>i}b(\mathbf{R}_{ij})\rangle=\dfrac12N_{\rho}\int_0^{\infinity}b(r)g(r)4\pi r^2\mathrm{d}r
\end{displaymath}
该公式对各向同性的液体成立，$\rho$是粒子密度。\textcolor{red}{Computer Simulation of Liquids by M. P. Allen and D. J. Tildesley}

将重新计算\textrm{MD}，计算前，将前次模拟得到的\textrm{CONTCAR}保存到\textrm{POSCAR}。

\textrm{INCAR}文件

\textrm{KPOINTS}文件

计算模拟参数与之前一样。

计算得到的对相关函数写入\textrm{PCDAT}文件，可以通过\textrm{Bash}脚本\textrm{pair\_correlation.sh}和\textrm{pair\_correlation.gp}通过\textrm{gnuplot}实现可视化。

前面的计算已经完成了\textrm{90fs}的模拟，因此将当前的e\textrm{PCDAT}为\textrm{PCDAT.90fs}。由当前最后的结构\textrm{(CONTCAR)}，开启重新计算，完成的计算是\ch{Si}熔融时长\textrm{180fs}的模拟，但模拟结果则是对\textrm{90fs}取的系综平均(计算控制参数\texttt{NSW}和\texttt{POTIM})。得到的\textrm{PCDAT}另存为\textrm{PCDAT.180fs}。然后就可以作图对比两个对相关函数。

从对相关函数结果看，前\textrm{90fs}，系统保持与初始结构接近，因为对相关函数中存在晶体才有的长程有序特征峰。有必要指出的是，当$r$超过超晶胞最短晶格一半的时候，$g(r)$对$r$的曲线没有意义。因此在本算例中$r$的取值应该是$\leqslant 5.4\AA$。

%长距离条件下的对相关函数
采用\textrm{py4vasp}可视化结构。

\textrm{XDATCAR}保存了每个时间步长的离子位置。粒子径迹是粒子在实际空间的运动轨迹。例如，下图中，根据灰色箭头，\textcolor{red}{cyan}色的原子离开中心的绿色盒子。

另一种表示\textrm{XDATCAR}文件是中心盒表示。在每个时间步长内，通过粒子位置加载周期边界条件，得到``封装的''离子径迹。在上述图示中，就是如红色虚线箭头示意的\textrm{cyan}原子被周期镜像取代。

将\textrm{CONTCAR}另存为\textrm{POSCAR}文件，并修改\textrm{INCAR}文件中的计算控制参数计算额外的\textrm{15fs}，注意\textrm{OUTCAR}中\texttt{NPACO}，\texttt{APACO}，\texttt{NBLOCK}和\texttt{KBLOCK}。这些计算控制参数可以改变最大距离(本算例中，对相关函数中该值为$5.4\mathrm{\AA}$)，根据最后的离子位置，重启计算。

因此\textrm{INCAR}中，设置\texttt{NSW}=5\texttt{APACO}=5.4，按照\textrm{Bash}脚本

用\textrm{Bash}脚本，通过\textrm{gnuplot}绘制新的对相关函数可视化图\textrm{pair\_correlation.png}。

最后\textrm{15fs}对应的是总模拟时间为\textrm{195fs}的对相关函数。可以看到这里的对相关函数比较差。

\subsubsection{分子的几何结构}
采用第一原理\textrm{MD}模拟\textrm{Langevin}等温条件下16个\ch{Ni}原子\textrm{NPT}系综下的几何结构的原子坐标。

\textrm{VASP}中通过计算控制参数\texttt{MDALGO}设定各种等温模拟，本算例将采用\textrm{Langevin}等温，因此\texttt{MDALGO}=3。\textrm{Lagevine}等温对粒子运动方程的贡献包括随机项和确定项，确定项来自\textrm{Nos\'e-Hoover}等温，而随机项来自\textrm{Andersen}等温(\textrm{MDALGO}=1)，\textcolor{red}{这就是以随机方式引入温度的影响的等温方式}。

在有很强的分子间相互作用力的体系，如弯曲振动模的模拟中，\textrm{Andersen}等温和\textrm{Langevin}等温\textcolor{red}{没有在各态遍历假设之上，引入新的内容}。实际上，\textcolor{red}{这两种等温方式，在平衡自由度上非常有效的}。相反，\textrm{Nose\'e-Hoover}等温在小体系或强直体系中，没有考虑各态遍历假设。例如，在单个丁烷模拟中就存在这个问题。实际上，在\textrm{VASP}模拟中，等温方式主要取决于热力学系综的选择。

在\textrm{MD}模拟中，有可能会想检查每个时间步长的原子位置变动。\textcolor{red}{原子位置的配置，主要归结为分子几何结构}。在\textrm{VASP}中，\textrm{ICONST}文件中确定的是体系内部的几何坐标。

在金刚石型\ch{Si}的初级原胞基础上，构建$2\times2\times2$的超晶胞，\textcolor{red}{在原有\textrm{Pymatgen}构建代码，增添下以下代码内并执行}

由此可以对照生成的\textrm{POSCAR}文件和参照\textrm{POSCAR.ref}。

输入控制文件\textrm{INCAR}

$\vec k$空间剖分网格文件\textrm{KPOINTS}

\textcolor{red}{\textrm{ICONST}文件}

在\textrm{INCAR}文件中，计算控制参数\texttt{ISYM}=0，计算过程中完全不考虑空间群对称性。\textrm{DFT}计算步中，计算精度控制参数\texttt{PREC}\textrm{=normal},\textrm{IVDW}=10考虑\textrm{Van der Waals}校正。此外的都是\textrm{DFT}重用的计算参数，为了节约存储空间，不写\textrm{WAVECAR}文件和\textrm{CHGCAR}文件。

\textrm{MD}计算的计算控制参数包括\texttt{IBRION}，\texttt{NSW}和\texttt{POTIM}。而\texttt{ISIF}=3和\texttt{MDALGO}=3，表示计算选择\textrm{Langevin}等温完成\textrm{NPT}系综模拟。

\textrm{ICONST}文件确定了待检查的分子结构。第一列确定了几何坐标的种类，如离子距离坐标原点的距离，或者是离子间的夹角。随后的整数，是\textrm{POSCAR}文件中相同顺序的原子序号。只有最后的整数是指定行为，即整数7，表示监察原子位置。

\textrm{KPOINTS}文件确定了一套均匀的$\vec k$空间网格剖分方案，包括$\Gamma$点，每个方向上两个格点。本该采用更大的计算晶胞，但计算开销太大。实际计算中，需要检查关注的物理量计算是否随晶胞结构和$\vec k$点收敛。因为这里只关心空间坐标的变化，所以这样的计算已经足够。

运行计算脚本，因为不仅是$\Gamma$点计算，主要要用\textrm{vasp\_std}。因为$vasp\_gam$只支持仅有$\Gamma$点的单个$\vec k$点，在\textrm{vasp\_gam}中开的数组是是实数类型的，而非复数类型的。而这里的均匀$\vec k$空间网格剖分点，不仅仅是$\Gamma$点。

在\textrm{REPORT}文件中，有每个\textrm{MD}步，有待检查原子坐标，例如，在第一个\textrm{MD}模拟步中的数据

采用\textrm{py4vasp}，提取检视的体积并作图。

\textrm{OUTCAR}文件中保存了模拟的总时间，

根据当前计算的结果，估计10000\textrm{MD}步的大致时间。从目前的结果看，10个\textrm{MD}步骤，大致是\textrm{74s}，所以10000\textrm{MD}步大约是{20}小时。

10000\textrm{MD}步的步数规模是一个典型的\textrm{MD}模拟得到合理结果的必要步数。因此，从这个算例可以看得出，第一原理\textrm{MD}模拟是非常吃计算开销，并且也非常耗时。所以很多体系采用第一原理\textrm{MD}模拟是不太可行的。

\subsection{机器学习力场}
在第一原理\textrm{MD}运行过程中，实时训练原子受力的力场。本算例是采用\textrm{Langevin}等温完成16个\ch{Si}原子体系\textrm{NPT}系综的第一原理\textrm{MD}模拟，在模拟中实现每个离子步实时修正力场。

第一原理\textrm{MD}模拟极其消耗计算开销。因为好的热力学平均，需要好的系综采样。换言之，模拟的时间极长。同时，每个第一原理\textrm{MD}步(离子步)的计算开销也很大，因为离子步包含的电子步采用\textrm{DFT}方法计算，属于完全的量子力学计算。前面的计算已经可以估算出\textrm{MD}的耗时程度。

在经典的\textrm{MD}理论中，并非采用\textrm{DFT}计算来得到每个原子上的作用力，而是通过力场\textrm{(Force Field)}来描述原子受力。力场也被称为原子间相互作用势\textrm{(Interatomic Potentials)}，可以通过实验观测或已有原子间作用力的经验总结得到，如\textrm{Van der Waals}相互作用，静电相互作用等等。而原子间相互作用势的优劣，取决于对体系相互作用的了解程度。显然，经典\textrm{MD}虽然依赖经验，但这样做大大降低了每个\textrm{MD}计算的开销，较长时间的\textrm{MD}模拟也变得可行。

源于第一原理的\textrm{MD}机器学习力场\textrm{(Machine learning Force Field)}既从第一原理层面，通过量子力学精确处理相互作用的电子能量(势能面)，精确刻画了粒子间相互作用的物理本质，又支持以较低的计算开销完成长的分子动力学模拟，可谓``鱼与熊掌兼得''的策略。

模拟结构\textrm{POSCAR}文件

计算控制参数\textrm{INCAR}文件

\textrm{ICONST}

$\vec k$空间网格剖分\textrm{KPOINTS}文件

s输入文件中关于\textrm{DFT}和\textrm{MD}计算的控制参数之外，机器学历力场的计算控制参数\texttt{ML\_LMLFF}=\textrm{T}，表示使用机器学习力场方法；而\texttt{ML\_ISTART}=0设定为采用从第一原理\textrm{MD}中训练新的力场的方式完成计算。\textcolor{red}{\textrm{RANDOM\_SEED}保证了本算例的可重复性。}实际上，\textrm{MD}模拟应有足够的时长，确保随机种子\textrm{(random seed)}是完全无关的。\textcolor{red}{不过自本例中，希望模拟时长尽可能短，并且，当训练的力场取决于随机种子时，原子需要面对的局域参考构型\textrm{(local reference configurations)}的数目也要加入学习-训练的过程。}局域参考构象受模拟时长的影响，因此也就会影响力场的性能。

\textcolor{red}{\textrm{ICONST}的意义与之前相同。}

\textrm{KPOINTS}文件确定的$\vec k$空间网格的均匀剖分方案。在机器学习方案中，如果学习的力是正确的，可以将小的晶胞上训练的力场，可以适用到大体系中。这就要求学习的力是随着$\vec k$空间网格剖分点收敛。后续会介绍力场在不同体系下的迁移。

为了让\textrm{VASP}将学习的力场作为输入数据读入，首先将\textrm{ML\_FFN}另存为\textrm{ML\_FF}；再重启\textrm{MD}模拟时，是使用力场或重新学习力场，就要特别注意参数\texttt{ML\_START}的正确设置。

\texttt{ML\_LOGFILE}文件保存了机器学习力场方法的全部日志数据。

参数输入样本误差\textrm{(In-sample error)}是训练集的平均误差，而输出样本误差\textrm{(out-of-samole error)}是考虑同样系综的随机新构象后出现的平均误差。对于训练得较理想的力场，输出样本误差应该达到极小。\textrm{Bayesian}误差本质是输出样本误差的估计，标志力场训练的理想程度。

采用\textrm{py4vasp}画出\textrm{Bayesian}误差和输入样本误差。

如果\textrm{Bayesian}误差图中存在大的尖峰，表明机器学习力场面对的局域原子构象偏离了\textrm{VASP}学习的数据。因此对照第一原理参比结果，用预测的力场得到的结果显然较差。

当训练力场的时候，如满足以下任一条件，\textrm{VASP}将由\textrm{MD}模拟得到结构作为机器学习的新的局域构象的参考态。
\begin{itemize}
	\item 如果尚没有机器学习力场存在，将晶胞结构的局域参考构象加入训练集，，用于构建力场；
	\item 在任何一个\textrm{MD}步，任何原子受力的\textrm{Bayesian}误差如果高于\texttt{ML\_CDOUB}$\times$\texttt{ML\_CTIFOR}确定的阈值，该结构的局域参考态构象也将被加入训练集，重新训练力场；
	\item 力场构造完毕，根据参数\texttt{ML\_NMDINT}\textrm{VASP}开始若干步\textrm{MD}模拟的预热步骤。在预热步中，\textrm{VASP}对于原子受力的\textrm{Bayesian}误差限制较小。当预热步骤完成后，如果原子受力的\textrm{Bayesian}误差是高于\texttt{ML\_CTIFOR}规定的阈值，\textcolor{red}{则该结构将被当作待考订结构}对待，当待考订结构数目达到\texttt{ML\_MCONF\_NEW}所规定的数目，就将新的局域参考构象会被加入到训练集。
\end{itemize}
必须指出的是，\textrm{Bayesian}误差是对输出采样数据的估计，并不是机器学习力场的质量标记。
、
\subsubsection{通过离子弛豫测试力场}
采用前面算例训练的机器学习力场计算16个\ch{Si}原子构成晶胞的若干个$E-V$曲线，并与\textrm{DFT}计算结果对比。

机器学习的算法目标之一，就是最小化输出采样数据误差，\textcolor{red}{即考虑训练集之外的数据后出现的误差。}只是训练过程中只考虑训练数据，因此只计算输入采样误差，这是训练集的平均误差。

\textrm{Bayesian}误差属于对输出采样误差的一种估计方案，\textcolor{red}{属于力场训练的优化驱动方向。}不过，相比于输出采样误差，\textrm{Bayesian}误差诉不上纯粹的力场质量的衡量标准。\textcolor{red}{为了完成前面力场的测试，可以采用其中一条、两条或三条策略}
\begin{itemize}
	\item 准备随机构象的独立测试数据集，通过\textrm{DFT}和机器学习力场的预测数据集分别计算两个结构得到的力、应力和能量差；
	\item 如有了第一原理参考数据和力场计算的结果，检查两者得到的物理性质的一致性。如弛豫的结构参数，声子，不同相的相对能量，弹性常数，缺陷组成等。
\end{itemize}
这个算例，对照的是后一种策略，对照弛豫的结构参数。

输入\textrm{INCAR}文件，保留前面的\textrm{ML\_FFN}

参考用的结构文件\textrm{POSCAR.ref}

\textrm{INCAR}文件

\textrm{KPOINTS}文件

在\textrm{INCAR}文件中确定为离子弛豫。

计算采用的\textrm{Bash}脚本。

运行上述脚本，可以采用以下脚本\textrm{volume\_energy.sh}，提取每次得到的体积和能量。该脚本的内容是：

运行脚本

对照机器学习力场\textrm{(Machine-Learned Force Field, MLFF)}和\textrm{DFT}得到的参照数据分别得到的总能和体积的$E-V$曲线图。\textrm{DFT}与机器学习训练训练的力场得到的数据是一样的。

\textrm{MLFF}结果与\textrm{DFT}参考数据非常接近。

\subsubsection{系综平均的晶格常数和晶胞体积}
使用机器学习力场使用\textrm{Langvin}常温在温度\textrm{400K}\textrm{NPT}系综平均下16个\ch{Si}原子的晶格常数和晶胞体积。

通常\textrm{DFT}计算的晶格常数和晶胞体积仅是$T=0\mathrm{K}$的基态性质。不过很多材料的结构特征会随着温度上升发生变化。更重要的是，实验上测量的晶格参数是热力学系综平均。采用前面训练的机器学习力场计算晶格常数和晶胞体积，这是统计力学可观测的物理量。力场训练的温度是$T=400\mathrm{K}$，晶格结构和晶胞体积计算也是在该温度下完成的。

初始结构文件\textrm{POSCAR.init}

初始输入文件\textrm{INCAR.init}

\textrm{KPOINTS}

\textrm{INCAR}文件中%第一原理计算没有特别的说明
确定了\textrm{MD}模拟，并使用机器学习力场。

\textrm{POSCAR}文件是\textrm{400K}等温的\ch{Si}晶体。

执行\textrm{Bash}脚本。

在\textrm{OUTCAR}文件中，查看最后10000离子步的计算耗时，估计总的耗时，并对照第一原理\textrm{MD}模拟的时间。

不到\textrm{10min}和\textrm{20h}的差别。

然后在\textrm{REPORT}文件中，检查每个\textrm{MD}离子步的几何结构坐标。通过\textrm{Bash}脚本\textrm{get\_lattice\_parameters.sh}提取结构参数，并对模拟时间取平均。

晶胞体积

晶格参数和平均晶格参数。

从计算结果看，得到的结果和理想对称性的晶胞的偏差，模拟的时间还应该增加。

\subsubsection{机器学习力场的迁移和热膨胀系数}
应用\textrm{400K}得到的机器学习力场，分别模拟\textrm{200K}和\textrm{300K}的16个\ch{Si}原子构成的晶格常数和晶胞体积，检验机器学习力场的迁移能力。分别作\textrm{200K}、\textrm{300K}和\textrm{400K}时的\textrm{E-V}曲线，确定体系的热膨胀系数。

训练好的机器学习力场，可以直接用于\textrm{MD}模拟而不再需要电子步的第一原理计算。如果在\textrm{MD}模拟中，出现偶尔的\textrm{Bayesian}误差增大，意味着模拟进入到了未知的相空间中。例如出现相变或结构变化，这种情况下有必要重新训练力场，以便模拟在新的相空间中的行为。

因为训练的数据到\texttt{ML\_ABN}文件中，为了继续训练力场，在\textrm{INCAR}中，设置计算控制参数\texttt{ML\_ISTART}=1,将\texttt{ML\_ABN}文件另存为\texttt{ML\_AB}。训练常常要分为很多轮，另外可能训练中会包含缺陷、表面或新的相空间内容。最终得到的力场参数保存在\textrm{ML\_FFN}文件中。

不过，当有些模拟控制参数，如温度，发生改变时，则未必需要重新训练力场。特别是如果全部局域参考构象足以覆盖全部系综，更没有此必要。特别是如果没有相变发生，这一点可以通过检验金刚石结构\ch{Si}来验证:~在训练得到机器学历力场的条件下，通过足够长的\textrm{MD}模拟，再把该力场平移到新的观测条件下，完成相同的模拟，考察两种条件下的晶胞体积的变化。

根据结果数据绘制的图像，构建$T=400\msthrm{T}$的力场所需要的全部局域构象，可以覆盖部分\textrm{200K}条件参数的部分局域构象。

$T=200\mathrm{K}$时的\textrm{POSCAR}结构

$T=300\mathrm{K}$时的初始结构\textrm{POSCAR.init}

$T=200\mathrm{K}$时的计算控制参数文件\textrm{INCAR.init}

$T=300\mathrm{K}$时的计算控制参数文件\textrm{INCAR.init}

\textrm{MD}模拟中改变的原子位置参数文件\textrm{ICONST}

$\vec k$空间网格剖分文件

在\textrm{ICONST}文件中，参数\texttt{LR}和\texttt{LV}分别是监控晶格常数和晶胞体积变化，\texttt{LA}监控的是晶格的角度。

\textrm{Bash}命令执行计算。

提取全部\textrm{MD}模拟的每个离子步的晶格矢量和晶胞体积，分别作系综平均。

编写\textrm{Python}脚本或使用现有\textrm{Bash}脚本\textrm{get\_lattice\_parameters.sh}% !表示注释\textrm{Bash}脚本
脚本的内容:

200\textrm{K}时，体积:
晶格常数:
平均晶格:

300\textrm{K}时，体积:
晶格常数:
平均晶格:

应用\textrm{Py4vasp}脚本，绘制晶胞体积-温度的变化曲线图，并与\textrm{400K}时的结果对照。

体积分别为：

在当前(200\textrm{K}-400\textrm{K})温度范围内，体系的晶胞体积$V$随温度$T$基本呈线性关系。因此可据此估计各向同性的膨胀系数:
\begin{displaymath}
	\alpha_V=\dfrac1{V}\dfrac{\mathrm{d}V}{\mathrm{d}T}
\end{displaymath}

换言之，由此可以得到体系体积对温度变化斜率随温度的关系。

热膨胀系数为\ch{Si}:~$3.46850\times10^{-6}$。

s实验中，线膨胀系数定义为:
\begin{displaymath}
	\alpha_L=\dfrac1L\dfrac{\mathrm{d}}L\mathrm{d}T$
\end{displaymath}
对于\ch{Si}的线膨胀系数实验为约$2.6\times10^{_6}$\upcite{JAP56-314_1884}。对于立方晶系，有等式
\begin{dieplaymath}
	\alpha_L=\alpha_V/3
\end{dieplaymath}
并据该式可以确定$\alpha_L$

并得到线性膨胀系数为$1.15617\times10^{-6}$。

同时，$\alpha_L$可以通过直接计算晶格矢量的长度得到。这里也可以利用立方晶体三个格矢等同，对计算值求平均，提升统计的精确程度。利用\textrm{Py4vasp}模块，写\textrm{Python}代码，作格矢对温度变化的曲线，得到$\alpha_L$，并与晶胞膨胀系数计算得到的值对比，可以判断哪个数值更合理。

\ch{Si}的线性膨胀系数为:~$1.38308\times10^{-6}~\mathrm{K}^{-1}$

\subsection{\ch{Cl-}对\ch{CH3Cl}的取代反应}
\subsubsection{自由能的缓慢增长模拟}
对采用缓慢增长\textrm{Slow-growth}模拟和热力学积分，对\ch{CH3Cl}被\ch{Cl-}的$\mathrm{SN}_2$亲核取代反应，计算作为反应坐标\textrm{(Reaction Coordinate)}函数的自由能。

计算化学反应的自由能变化曲线，可以得到反应的活化自由能\textrm{(Activation Free Energy)}。化学反应伴随了稳定的反应物变到过渡态\textrm{Transition State}。得到自由能曲线的办法称为热力学积分\textrm{(Thermodynamic Integration)}。

通过计算化学反应过程中的自由能曲线，可以得到反应的活化自由能\textrm{(activation free energy)}。活化自由能就是从稳定的反应物形态转化为产物形态所需要做的功。在\textrm{MD}模拟中，得到反应的自由能曲线的方法称为热力学积分，是沿着反应路径\textrm{(Reaction Path)}对自由能梯度进行积分计算。反应路径(也成为反应通道)是联结化学反应的起始和终了的热力学态的能量面，习惯上，将反应通道简化为反应物向产物变化的能量变化曲线，称为反应坐标\textrm{(Reaction Coordinate)}。

缓慢增长模拟采用第一原理\textrm{MD}模拟来确定反应物到产物的转变过程及机理，系统由参与反应的物质构成，假设体系渐次被活化的单步能量只有旋转自由度能级，确保化学反应可以非常缓慢地发生，并能满足反应过程中每个时刻都保持可逆的要求。根据热力学要求，这样才满足对体系做的功等于反应的自由能。

与之相反的，同样存在快速切换\textrm{Fast-switching}模拟。这种情况下采用多次快速完成对过渡态的纯静态自由能计算实现的。这样做导致对反应体系做的功是不可逆的。通过\textrm{Jarzynski}恒等式\textrm{(Jarzynski equality)}\footnote{\textrm{Jarzynski}恒等式是非平衡统计力学中连接非平衡做功过程与平衡态自由能差的核心等式。}，这些不可逆功可变成可逆反应的自由能。

另外还有一种方法是通过分析势能面\textrm{(Potential Energy Surface, PES)}在稳定点附近的局域\textrm{topology}结构。这种静态方法只需要对势能面上一些稳定点的旋转、振动能级进行模拟和分析，因此计算开销也小得多。不过，该方法尚未考虑出现在过渡态中的被微小的转动、振动模激发的反应，所以只能用于定性估计类似的反应路径：

\textcolor{red}{\huge 有机化学反应!!!!}

这是一个$\textrm{SN}_2$反应。氯代甲烷\textrm{(Chloromethane)}\ch{CH3Cl}和\ch{Cl-}碰撞，氯代甲烷成为亲核体，与\ch{Cl-}成键；逐渐形成松散的，但可原位观测的稳定中间体态；最终，离子态的\ch{Cl-}取代氯代甲烷中原来的\ch{Cl}，取代反应完成。$\mathrm{SN}_2$反应是对称的，一个\ch{Cl-}从一侧进攻，而另一个\ch{Cl-}从另一侧离开，其间伴随了\ch{CH3}结构的翻转。因此反应的自由能曲线也将是对称的。

初始结构文件\textrm{POSCAR.init}

初始计算控制文件\textrm{INCAR.init}

模拟反应的计算控制参数文件\textrm{INCAR.refit}

\textrm{ICONST}文件

\texttt{ML\_AB}文件中保存的是已经完成的\ch{CH3Cl}和\ch{Cl-}阴离子的第一原理数据，包括:~晶格矢量，原子位置，能量，力和应力张量。

\textrm{INCAR.refit}文件中与\textrm{MLFF}方法有关的控制计算参数为
\begin{itemize}
	\item \texttt{ML\_LMLFF}=\textrm{T}，表示开启并使用\textrm{MLFF}；
	\item \texttt{ML\_MODE}=\textrm{refit}表示允许在现有\texttt{ML\_AB}文件基础上产生力场文件。
\end{itemize}

\texttt{SYSTEM}只是描述允许计算体系，允许是任意字符串。\texttt{ISYM}在\textrm{MD}模拟中通常设置为关闭对称性。\texttt{POMASS}指定了参与反应的原子/离子质量，默认为\textrm{POTCAR}中的元素质量。如果设定了该值，则不使用\textrm{POTCAR}中的默认值。之所以在化学反应\textrm{MD}模拟中改变某些元素的质量，比如这里对于三个\ch{H}元素，质量比\ch{H}的\textrm{POTCAR}中的质量重，因为如果适当减缓\ch{C-H}键的振动频率，可以通过\texttt{POTIM}使用更大的时间步长，而不会损失反应过程中的动力学信息。

这种模拟中轻微改变元素质量，与同位素效应\textrm{(Isotope effect)}不同，并不改变第一原理\textrm{MD}模拟的力场。
%\footnote{所谓同位素效应，会改变实际反应速率的，本质上是离子自由度的量子效应。必须考虑从量子力学角度考虑离子的处理。原则上，机器学习力场是可以重现第一原理\textrm{MD}的，两者是等效的。除非训练的力场不够充分，否则两者的结果并没有区别。在第一原理\textrm{MD}模拟中，离子是作为经典粒子对待的，而每个离子所受的力则是通过量子力学方法计算得到的。最后，在经典力学中，热力学平均的方程中，离子质量也会彼此抵消，因此无法在\textrm{MD}中观察到同位素效应。}

%\textrm{INCAR}文件中计算控制参数\texttt{KSPACING}确定了$\vec k$点网格密度，将该值设置为大得离谱的数值，本质上只包含$\Gamma$点。当然可以通过\textrm{KPOINTS}文件确定$\vec k$空间剖分网格。

\texttt{IBRION}，\texttt{NSW}和\texttt{POTIM}控制了模拟中每个离子步离子位置的更新方式。\texttt{MDALGO}，\texttt{ANDERSEN\_PROB}，\texttt{TEBEG}和\texttt{ISIF}共同确定了热力学系综。\textrm{Andersen}等温采用随机碰撞来模拟热浴行为。

\textrm{ICONST}文件给出具体的反应坐标中原子/离子的成键关系。在本算例中，\textrm{POSCAR}中，原子1是\ch{C}，原子5与6是\ch{Cl-}，\textrm{ICONST}文件中确定的反应坐标，就是有两种不同的\ch{C-Cl}键的键长。这意味着，如果在\ch{CH3Cl}中\chP{C-Cl}键的键长为\textrm{1.8\AA}，而\ch{C}与被取代的\ch{Cl-}离子的距离为\textrm{3.0\AA}，则反应坐标为的值为\textrm{1.8\AA-3.0\AA=-1.2\AA}。在过渡态\textrm{Transition State}，两个\ch{Cl-}离子距离\ch{C}具有相同的距离，因此反应坐标为0.

\textcolor{red}{在完成一个化学反应模拟的过程中，设计合适的反应坐标是各类反应动力学模拟最重要的任务。}

除了\textrm{ICONST}中规定的原子成键关系，\textrm{INCAR}文件中\texttt{INCREM}参数限制了发生的反应沿反应通道演化的速度，需要说明的是，该速度指的是每个离子步长内的速度，而非单位时间的速度。

设置\texttt{LBLUEOUT}参数会在\textrm{REPORT}文件中输出计算的自由能梯度数据。

运行\textrm{Bash}脚本

将第一原理计算的能量、力和应力保存到\texttt{ML\_AB}文件中，生成机器学习力场。

将产生的机器学习力场由输出文件\textrm{ML\_FFN}复制到\textrm{ML\_FF}文件中，执行\textrm{Bash}脚本。

将前面运行的\textrm{CONTCAR}作为反应物初始结构\textrm{POSCAR}；并用首次计算最终的\texttt{RANDOM\_SEED}作为第二次计算的原子初始速度。

这两步确保后续计算在第一次计算基础上完成后续计算。

%执行3000\textrm{MD}步，
关于热力学积分:~一旦需要在特定的反应坐标$\xi^{\prime}$计算自由能梯度$\dfrac{\partial A}{\partial\xi}$，可以通过对反应路径积分，计算反应中任意两个态\textrm{A}和\textrm{B}的自由能差分$\Delta_{\mathrm{A}\rightarrow\mathrm{B}}$
\begin{displaymath}
	\Delta A_{\mathrm{A}\rightarrow\mathrm{B}}=\int_{\xi_{\mathrm{A}}}^{\xi_\mathrm{B}}\left.\dfrac{\partial A}{\partial\xi}\right|_{\xi^{\prime}}\mathrm{d}\xi^{\prime}=\Delta A(\xi_{\mathrm{B}})-\Delta A(\xi_{\mathrm{A}})
\end{displaymath}

热力学积分的主要困难在于在每个$\xi^{\prime}$点都能有足够多好的自由能梯度采样$\dfrac{\partial A}{\partial\xi}$。采用缓慢增长方法，反应坐标$\xi^{\prime}$都随着时间步长(\texttt{POTIM}设定)$\mathrm{d}\tau$改变，反应速率(由\texttt{INCREM}设定)为$\dot{\xi}^{\prime}$。因此上式在缓慢增长近似下，引入
\begin{displaymath}
	\mathrm{d}\xi^{\prime}=\dot{\xi}^{\prime}\mathrm{d}\tau
\end{displaymath}
该方法不适合转变速率和时间步长都不是足够小的状态。\textrm{MD}模拟的质量可以通过反应过程的可逆程度来检验。\textcolor{red}{如果自由能曲线是由$\Dalta A(\xi)$绘制得到是迟滞曲线，表明转换速率和时间步长都不是足够小。}

实际上，缓慢增长模拟只是用来快速了解反应自由能曲线的。为了完成一个高质量的热力学积分，\textrm{VASP}中还有其余更合适的方法，比如``蓝月亮''\textrm{(blue-moon)}方法\footnote{蓝月亮方法是一种基于约束分子动力学的自由能计算方法，用于沿指定反应坐标计算化学反应的平均力势(\textrm{PMF}，即自由能剖面)}，可以更精确地控制收敛，但计算成本的开销更大。

完成1500步\textrm{MD}模拟，会出现一个目录\textbf{run1}，打开\textbf{run1}\textrm{/REPORT}文件。在每个\textrm{MD}步，有以下数据：

\textcolor{red}{作为缓慢增大模拟的后处理，将反应坐标起始点作为${\xi^{\prime}}$态，在$\xi^{\prime}$的自由能梯度$\dfrac{\partial A}{\partial\xi}$是曲线起始点的梯度。}由此可以开启积分。

采用\textrm{Py4vasp}实现\textrm{MD}径迹的可视化。

继续完成1250步\textrm{MD}模拟，又有一个目录\textbf{run2}生成，完成计算后可以根据\textrm{Bash}脚本\textrm{free-energy\_gradient.sh}提取自由能梯度：

执行脚本产生\textrm{free-energy\_gradient.dat}

采用\textrm{Python}脚本完成热力学积分，绘制反应自由能作为反应坐标函数的曲线。

生成的自由能作图中，$\xi<0$，模拟的是反应物，$\xi>0$，模拟的是产物。$\xi=0$对应于过渡态。

因为反应是$\mathrm{SN}_2$是对称取代反应，所以自由能曲线是对称的。本次模拟得到的反应通道偏离理想对称点，归因于时间步长(\texttt{POTIM})或者变换速率(\texttt{INCREM})太大。
从\textrm{MD}径迹中每个态里找到代表性的构象。

依次检验起始步作为反应物，近似750步左右，作为反应过渡态，到1499步作为产物。
\subsubsection{反应物态的概率分布}
完成反应物\ch{CH3Cl}和\ch{Cl-}离子的自由\textrm{MD}模拟，并计算反应坐标的概率密度。

如果\ch{CH3Cl}和\ch{Cl-}离子间的反应未经活化，反应不可能自发进行。该结论可以通过反应物系综内计算对反应坐标的概率密度
\begin{displaymath}
	P(\xi^{\prime})=\langle\delta(\xi^{\prime}-\xi(\mathbf{R}))\rangle_{\mathbf{R}^{\prime}}, \xi^{\prime}<0
\end{displaymath}
得到验证。

通过\textrm{Bash}脚本，建立\textrm{ML\_FF}文件的软链接。

初始结构文件\textrm{POSCAR.init}

初始计算控制文件\textrm{INCAR.init}

\textrm{ICONST}

力场参数文件\textrm{ML\_FF}

这里的\textrm{INCAR}文件中，没有设置计算控制参数\texttt{LBLUEOUT}和\texttt{INCREM}，因此自由能梯度不会写入到\textrm{REPORT}文件，\textrm{MD}计算也不会强制沿着反应通道进行。也就是说，这是基于机器学习力场的自由非约束的\textrm{MD}模拟。

执行计算的\textrm{Bash}脚本

起始10000步\textrm{MD}。为计算反应速率常数\textrm{(Reaction rate constant)}和实验可检测的活化能。

化学反应的反应速率常数由反应物\textrm{R}和产物\textrm{P}共同确定。实验上可测量的反应速率常数，满足\textrm{Eyring}方程
\begin{displaymath}
	k_{\mathrm{R}\rightarrow\mathrm{P}}=\dfrac{k_{\mathrm{B}}T}h\mathrm{exp}\bigg(-\dfrac{\Delta A^{\ddagger}}{k_{\mathrm{B}}T}\bigg)
\end{displaymath}
这里$\Delta A^{\ddagger}$是唯相的活化能，根据\textrm{Arrhenius}定理，可以通过实验中反应速率常数与温度$T$的数据得到，$h$是\textrm{Planck}常数，$k_{\mathrm{B}}$是\textrm{Boltzmanna}常数，$T$是体系的温度。在\textrm{MD}模拟中，反应速率常数可以定量地表示为
\begin{displaymath}
	k_{\mathrm{R}\rightarrow\mathrm{P}}=\dfrac{\langle|\dot\xi|\rangle}2P(\xi_{\mathrm{ref},\mathbf{R}})\mathrm{exp}\bigg(-\dfrac{\Delta A_{\xi_{\mathrm{ref},\mathbf{R}\rightarrow\xi^{\ast}}}}{k_{\mathrm{B}}T}\bigg)
\end{displaymath}
这里$\Delta A_{\xi_{\mathrm{ref},\mathbf{R}\rightarrow\xi^{\ast}}}$是反应物在参考态$\xi_{\mathrm{ref},\mathbf{R}}$和过渡态$\xi_{\xi^{\ast}}$的自由能之差。这个指可以通过自由能曲线的计算直接得到。

只要体系还是反应物，选择的参考态是任意的，反应速率与所选参考态无关。因为反应速率正比于在参考态找到当前系统的概率；但反应速率与过渡态的广义速率的绝对值的热力学统计平均$\langle|\dot{\xi}^{\ast}|\rangle$有关。后面将会讨论如何计算$\langle|\dot{\xi}^{\ast}|\rangle$。

由此得到实验测量的唯相活化能与\textrm{MD}模拟的活化能的关系为
\begin{displaymath}
	\Delta A^{\ddagger}=\Delta A_{\xi_{\mathrm{ref},\mathbf{R}\rightarrow\xi^{\ast}}}-k_{\mathrm{B}}T\ln\bigg(\dfrac{h}{k_{\mathrm{B}}T}\dfrac{\langle|\dot{\xi}^{\ast}(0)|\rangle}2P(\xi_{\mathrm{ref},\mathbf{R}})\bingg)
\end{displaymath}
理论值与实验值对比时，有个常见的错误就是遗漏了上式中的第二项。

完成5000步\textrm{MD}模拟，会出现一个名为\textbf{run1}的计算目录。因为没有设置\texttt{LBLUEOUT}，因此\textrm{REPORT}文件中没有\textrm{MD}模拟\textcolor{red}{关于反应坐标约束\texttt{cc>}引领的数据。}取而代之的是\texttt{mc>}引领的记录反应坐标的数据。

采用\textrm{Py4vasp}可以可视化\textrm{MD}模拟中的原子径迹。

计算完成后，会有第二个目录\textbf{run2}出现。通过\textrm{Python}脚本，可以从\textrm{REPORT}文件中提取反应坐标并绘图。

不难看出，反应不会自发发生。在产物区$\xi^{\prime}>0.1\mathrm{\AA}$的概率密度为零也证明了这一点。因为反应过程的自由能能垒远高于$k_{\mathrm{B}}T$，所以反应是吸热过程，在可观测时间长度内，直接的\textrm{MD}模拟，也几乎不可能看到一次翻越势垒。考虑到取值范围$[-1.5\AA,-1.49\AA]$处，$P(\xi^{\prime})\approx1.54$，所以反应和活化过程都是稀有事件\textrm{(Rare events)}。
\subsubsection{过渡态的平均速率}
计算\ch{CH3Cl}与\ch{Cl-}离子形成过渡态的平均广义速度，反应速率和唯相活化能。

反应速率常数可以写作
\begin{displaymath}
	k_{\mathrm{R}\rightarrow\mathrm{P}}=\dfrac{\langle|\dot{\xi}^{\ast}|}2P(\xi_{\mathrm{ref},\mathbf{R}})\mathrm{exp}\bigg(-\dfrac{\Delta A_{\xi_{\mathrm{ref},\mathbf{R}\rightarrow\xi^{\ast}}}}{k_{\mathrm{B}}T}\bigg)
\end{displaymath}
因此唯相的活化自由能$\Delta A^{\ddagger}$和可逆做功$\Delta A_{\xi_{\mathrm{ref},\mathbf{R}\rightarrow\xi^{\ast}}}$的关系为
\begin{displaymath}
	\Delta A^{\ddagger}=\Delta A_{\xi_{\mathrm{ref},\mathbf{R}\rightarrow\xi^{\ast}}}-k_{\mathrm{B}}T\ln\bigg(\dfrac{h}{k_{\mathrm{B}T}}\dfrac{\langle|\dot{\xi}^{\ast}|\rangle}2P(\xi_{mathrm{ref},\mathbf{R}})\bigg)
\end{displaymath}
因此，根据上述计算，可以通过\textrm{MD}模拟得到$k_{\mathrm{R}\rightarrow\mathrm{P}}$和$\Delta A^{\ddagger}$，唯一未知的量是过渡态的广义速度的热力学平均，其定义为
\begin{displaymath}
	\dot{\xi}=\sum_{i=1}^N\sum_{\mu=x,y,z}\dfrac{\partial\xi}{\partial\mathrm{R}_{i,\mu}}\dot{\mathrm{R}}_{i,\mu}
\end{displaymath}
这里$\mathrm{R}_{i,\mu}$表示原子$i$的\textrm{Cartesian}坐标。

\textcolor{red}{与前面的计算类似}，执行\textrm{Bash}脚本

初始结构文件\textrm{POSCAR.init}

初始计算控制参数文件\textrm{INCAR.init}

\textrm{ICONST}

机器学习力场\textrm{ML\_FF}。

\textrm{POSCAR}确定了过渡态位置的各原子的结构关系

\textrm{INCAR}文件中，未设置\textrm{INCREM}。\textrm{INCAR}与\textrm{ICONST}文件设置相结合，就可以完成\textcolor{red}{受限的}过渡态的\textrm{MD}模拟。

执行\textrm{Bash}脚本运行

首先执行6000步\textrm{MD}模拟。根据前面的计算提取$\Delta A_{\xi_{\mathrm{ref},\mathbf{R}}\rightarrow\xi^{\ast}}$的数值。比如可通过以下\textrm{Python}代码提取$\xi_{\mathrm{ref},\mathbf{R}}=-1.5\AA$相关的数据。

根据上例，得到$P(\xi_{\mathrm{ref},\mathbf{R}})$，即$[-1.5\AA,-1.49\AA]$处，$P(\xi^{\prime})\approx1.54$。

根据\textrm{SHAKE}算法，广义速度计算方式如下

由质量求逆的度规张量\textrm{(metric tensor)}
\begin{displaymath}
	Z=\sum_{i=1}^N\dfrac1{m_i}\sum_{\mu=x,y,z}\bigg(\dfrac{\partial\xi}{\partial\mathrm{R}_{i,\mu}}\bigg)^2
\end{displaymath}
因此广义速度可以表示为
\begin{displaymath}
	\langle|\dot{\xi}^{\ast}|\rangle=\sqrt{\dfrac{2k_{\mathrm{B}}T}{\pi}}\dfrac1{\langle Z^{-1/2}\rangle_{\xi^{\ast}}}
\end{displaymath}
\textcolor{red}{这些物理量的计算脚本如下!!!}

采用\textrm{Py4vasp}模块，按照计算得到的广义速度对\textrm{MD}模拟步数绘图，并判断结果是否收敛。

\textcolor{red}{作为例子，这个计算的收敛足够。当读者自己计算时，一般需要更精确地观察所以计算量对\textrm{MD}步数、\textrm{fs}单位的单位步长以及其余重要的参量有更好的收敛结果。例如，可以根据\textrm{Mann-Kendall}趋势测试，确保数据并没有出现特定的趋势，即平均值不会出现漂移。}

计算反应速率常数$k_{\mathrm{R}\rightarrow\mathrm{P}}$的\textrm{Python}代码如下:

计算得到的反应速率常数

活化能

这些数值可以与文献中第一原理\textrm{MD}模拟得到的体系数值对比。采用``蓝月亮''方法得到的预测实验活化能为0.4182~\textrm{eV}。注意，由于$\Delta A^{\ddagger}$的微小偏差会对$k_{\mathrm{R}\rightarrow\mathrm{P}}$有很大的影响，所以这两类数据一般不能直接对比。

\section{机器学习力场}
结合第一原理\textrm{MD}的机器学习力场可以节约模拟时间，并依然从第一原理更深刻地认知物理。这里介绍\textrm{MLFF}中的误差分析与超参数优化。

\subsection{训练集与测试集的误差分析}
对于给定的参比数据，计算\textrm{MLFF}模型的训练集和测试集误差。

\textrm{MLFF}模型的误差评估了训练集模型的精度，以及模型外推到训练集以外数据的可靠性。机器学习方法中，需要区分两种误差，训练集误差和测试集误差。训练集误差是用训练集评估的，训练集的误差是用\textrm{DFT}计算值与\textrm{MLFF}预测值差值来计算的。

测试误差则是通过未参与训练过程的外部测试集数据来评估的。一般认为采用\textrm{MLFF}计算测试集中的采样数据，在同等条件下，可以得到比较好的数据。\textcolor{red}{例如，即使一开始的训练是在原子数目少一点的体系上构建的，测试集中的结构应该拥有相同数目原子结构的描述能力。而且测试集要能得到描述对象的某些热力学相。}通过这种方式构建的测试集误差，能够给出由训练集推演的\textrm{MLFF}模型的结果到底有多好。

一般地，对比训练集误差与测试集误差，一般是三种结果：
\begin{itemize}
	\item 如果训练集的误差较小而测试集的误差较大，说明存在过拟合\textrm{(Overfitting)}，这种现象说明通过机器学习得到力场，不适合处理训练集以外的结构，外推性差。必须通过增加训练集或调节超参数来提升力场的预测能力；
	\item 训练集误差和测试集误差大致相当，表明对于应用体系，只要没有引起明显的误差，那么该力场就是适用的；
	\item 训练误差加大，相反测试误差较小，意味着测试集出现了偏差，不具备普适性。
\end{itemize}
\textcolor{red}{关于机器学习方法的误差分析，参见文献\cite{DOI:10.1063/5.0160326}，其中有关于\textrm{MLFF}的专门讨论。}

本算例中，将会根据保存在\texttt{ML\_AB}文件中的已有的第一原理数据重新拟合\textrm{MLFF}模型。通过\txtrm{ML\_AB}中的数据，分析训练集误差和测试集误差，也就是评估得到的力场。

%根据\textrm{VASP}产生的第一原理数据和\texttt{ML\_AB}文件，拟合\textrm{MLFF}模型。根据之前的计算结果，而不重新启动\textrm{DFT}计算。

计算控制文件\textrm{INCAR.refit}

计算控制文件\textrm{INCAR.run}

结构文件\textrm{POSCAR}

$\vec k$空间剖分网格文件\textrm{KPOINTS}

\texttt{ML\_AB}之前完成的\ch{LiH}体系的第一原理计算数据，包括晶格矢量，原子位置，能量，力和应力。

\subsubsection{重新拟合的输入说明}
重新拟合的输入控制文件\textrm{INCAR.refit}只有计算控制参数\texttt{ML\_LMLFF}确定机器学习模式，\texttt{ML\_MODE}选择拟合方式。一般地，重新拟合需要在实时训练完成\texttt{ML\_AB}文件之后才能开始，这样得到的\textrm{MLFF}更精确，因为此时的模型权重是通过线性回归\textrm{Linear Regression}确定的。%
采用重新拟合后的\textrm{MLFF}，\textrm{VASP}在计算时会更快(\texttt{ML\_MODE}=\textrm{run})，因为重新拟合会引入更多的功能算法。

\textrm{POSCAR}文件不会对重新拟合产生影响，因为重新拟合是根据第一原理计算的\texttt{ML\_AB}文件的数据实现的，而\textrm{POSCAR}文件仅是用于预测步，因此\textrm{OUCAR}中的能量和力会受到其影响。但在分析\textrm{MLFF}的时候不会注意到这些结果。

\textrm{KPOINTS}文件中仅含$\Gamma$点

重新拟合的时候，不会涉及\textrm{POTCAR}文件。机器学习算法会根据\texttt{ML\_AB}文件确定原子，这些原子可以和\textrm{POTCAR}中的原子无关。

\subsubsection{误差分析输入说明}
测试集的文件位置

包括\textrm{POSCAR}结构文件和\textrm{DFT}计算数据

本算例的测试集包括50个独立的结构，实际应用中，考虑到最终力场的应用场景，测试集应该更大些，比如差不多500个结构。

使用\textrmP{Bash}脚本确定训练集误差

重新拟合后，\textrm{VASP}会将力场写入\textrm{ML\_FFN}文件，而训练集误差则写入\texttt{ML\_LOGFILE}文件中。采用\textrm{Bash}命令检查\texttt{ML\_LOGFILE}文件，

\textrm{grep ERR ML\_LOGFILE}

进行误差分析：

下表给出三种力场的测试集误差
最后一行有5列数据，第三个是能量的训练集误差，单位是\textrm{eV}；第四个是力误差，单位是\textrm{eV/AA}；最后一个是应力张量的误差，单位是\textrm{kbar}。

下面为了分析测试误差，并与训练误差对比。将得到的\textrm{MLFF}也能够用到测试集中，应用\textrm{Bash}脚本，分别计算测试数据集的能量、力和应力。

相应的\textrm{Bash}脚本为\textrm{test\_set\_analysis.sh}

计算完成后，将通过\textrm{Py4vasp}模块执行\textrm{Python}命令，对比\textrm{DFT}数据和\textrm{MLFF}数据的误差。并将每种构型的测试集误差都绘图给出。

为计算总的测试误差对测试集中构型的平均值，执行\textrm{Python}脚本

由此可以得到的结论:~相比于测试集误差，力的训练集误差更小一些。不过因为两类误差都小于\textrm{10~meV/AA}，\textrm{MLFF}看起来还比较可靠。测试集误差与训练集误差总体上在同一数量级，所以\textrm{MLFF}可以用于训练集以外的结构计算。

\subsection{通过超参数优化精度}
所谓超参数，就是在\textrm{MLFF}模型中由用户定义，但无法通过训练\textrm{MLFF}来优化的参数。对超参数的优化，可以在改善\textrm{MLFF}的性能的同时，提升模拟精度。超参数包括截断半径\texttt{ML\_RCUT1}和\texttt{ML\_RCUT2}，它们描述的是计算的局域原子环境。%有关超参数及其优化，建议参考文献\cite{Hyperparameter optimization}。

根据经验，以下超参数对于提升\textrm{MLFF}的性能和精度最重要

\textcolor{red}{这里只是给出\textrm{VASP}中的少量超参数。}
\begin{table}[h!]
    \centering
    \begin{tabular*}{\textwidth}{@{\extracolsep{\fill}}c@{\extracolsep{\fill}}|@{\extracolsep{\fill}}c}
        \toprule
	超参数 &参数功能 &参数分类\\
\midrule
\texttt{ML\_RCUT1} & 双粒子间截断半径 &描述符\textrm{(descriptor)}\\
\texttt{ML\_RCUT2} & 三粒子间截断半径 &描述符\textrm{(descriptor)}\\
\texttt{ML\_WTIFOR} & 原子受力权重因子 &拟合参数\textrm{(fitting)}\\
\texttt{ML\_WTOTEN} & 总能量权重因子 &拟合参数\textrm{(fitting)}\\
\texttt{ML\_WTSIF} & 应力张量权重因子 &拟合参数\textrm{(fitting)}\\
\texttt{ML\_EPS\_LOW} & 局域参考构型稀疏因子 &稀疏性约束参数\textrm{(sparse)}\\
\texttt{ML\_EPS\_LOW} & 描述符型稀疏因子 &稀疏性约束参数\textrm{(sparse)}\\
        \bottomrule
    \end{tabular*}
    \caption{\textrm{INCAR}常用参数的简要描述}
\label{Tab:Hyperparameter-Tag}
\end{table}

这个例子主要是优化\texttt{ML\_RCUT1}提升力场的性能。

交叉验证\textrm{(Cross-Validation)}是判断超参数质量的一种方法。作为一个准备步骤，检查训练集误差没有远小于测试集误差，以防止过拟合。然后系统地改变超参数，寻找最优值。如果只是随机地变化超参数，不可能提升\textrm{MLFF}的质量。最后再对比不同超参数下训练集误差和测试集误差。

以优化\texttt{ML\_RCUT1}为例，对一系列的\texttt{ML\_RCUT1}的值，计算训练集误差(存于\texttt{ML\_LOGFILE文件})和力场(存于\textrm{ML\_FF}文件)。通过\textrm{Py4vasp}可以对\textrm{ML\_FF}文件中测试集进行误差分析。在参数改变过程中，只考虑训练集误差就够了。然后只需要基于\textrm{MLFF}的训练集的最小误差来分析测试集误差。

\textrm{INCAR.refit}

\textrm{INCAR.run}

\textrm{POSCAR}

\texttt{ML\_AB}集结了前面计算的\ch{LiH}的第一原理数据:晶格矢量，原子位置，能量，力和应力。

这里最重要的输入文件是\texttt{ML\_AB}和\textrm{INCAR}文件，而\textrm{VASP}计算必须的\textrm{POSCAR}和\textrm{POTCAR}不会影响超参数的优化。为了不使用\textrm{KPOINTS}，这里在\textrm{INCAR}文件中使用了计算控制参数\texttt{KSPACING}。无论何种方式剖分$\vec k$空间网格，倒空间晶格与重构力场无涉。

这里的\texttt{ML\_AB}文件来自前面的计算，这里针对双体描述符\texttt{ML\_RCUT1}，\texttt{ML\_AB}保存的第一原理计算数据将会用于\texttt{ML\_RCUT1}取不同值时的重新拟合过程。该任务可以通过执行\textrm{Bash}脚本完成。

\textrm{run\_hyperparam\_scan.sh}

执行该脚本，分别针对各个\texttt{ML\_RCUT1}的计算。

计算完毕后，再绘制训练集误差与截断半径的关系图。

\textrm{Bash}脚本为\textrm{extra\_training\_error.sh}

由此得到\textrm{training\_error.dat}数据文件，可以通过\textrm{Python}脚本绘制数据图，除了计算的数据，还有一个包含更多数据点的\textrm{training\_error\_store.dat}文件。这里作为算例，为解决计算开销计，只给出三个数据点。

绘制的图中含有三类不同的误差。能量误差，力的误差和应力误差。途中红色部分是这个算例中的数据点。原则上，如果三个误差都是足够小，就是好的结果。但也有可能用户只关注其中三种误差中的某一个的情况。例如，对于缺陷形成能\textrm{(defect-formation energies)}，\textrm{MLFF}模型就应该最小化能量误差；相反，声子频率计算会要求\textrm{MLFF}最小化力的误差。为了突出某一种误差，可以在拟合时引入权重系数。拟合中增加能量误差的重要性，可以调节计算控制参数\texttt{ML\_WTOTEN}。类似地，力的重要性，通过计算控制参数\texttt{ML\_WTIFOR}增加其权重，而应力误差的权重则是\texttt{ML\_WTSIF}。如果是关注体积预测的精确程度，就会关注应力张量的误差。

根据计算结果绘图，可以看出\texttt{ML\_RCUT1}=13是最精确的结果。因此保留\textbf{ML\_ABN\_13.0}和\textbf{ML\_FFN\_13.0}，留作进一步分析。实际上，继续执行\textrm{Bash}脚本，可以重复测试集分析的步骤

其中\textrm{test\_set\_analysis.sh}脚本为:

原则上，为了得到最好的\textrm{MLFF}，需要\textrm{VASP}要对每一个可能的超参数重复上述优化流程。但实际操作时，不可能优化全部超参数。使用\textrm{VASP}时，有默认的参数选项，这都是在一系列体相和分子数据集上优化过。\textcolor{red}{根据力场的计算预测能力，考虑上表中某些超参数优化是值得的。}

\subsection{优化\rm{MLFF}的性能}
根据现有的\texttt{ML\_AB}文件，优化力场性能。

有效的稀疏化提升力场的性能。例如，\textrm{VASPj通过计算控制参数\texttt{ML\_LSPARSDES}，采用\textrm{CUR}算法降低\textrm{MLFF}模型的角度描述符的数量。该协方差矩阵描述符$A$可分解为三个矩阵$C$，$U$和$R$，这里$C$是矩阵$A$的少数几行，$R$则是$A$的少数几列，而$U$则是一个小的幺正矩阵(酉矩阵，\textrm{unitary matrix})，它使得$CUR$的乘积接近矩阵$A$。%最初的\textrm{CUR}算法是用于有效筛选矩阵$A$中的重要行，这里\textrm{CUR}算法用于搜索不重要的角度描述符并筛除。通过计算控制参数\texttt{ML\_RDES\_SPARSDENS}和杠杆值\textrm{(leverage scores)}\footnote{杠杆值或杠杆得分衡量矩阵中每个数据点对拟合结果的影响力，常用于异常点检测和大规模矩阵计算的随机采样。}来考量。

关于\textrm{CUR}算法的详细内容，参见文献\cite{CS106-697_2009, JCP152-234102_20220}。

筛除描述符提升力场的计算能力的代价是降低计算精度。为了平衡精度和计算能力，可以将测试集精度为$y$轴，代码的时间开销为$x$轴作图，求出该多参量问题的\textrm{Pareto}前沿\textrm{(Pareto Front)}。根据得到到\textrm{Pareto}前沿，选定优化的\texttt{ML\_RDES\_SPARSDES}值。为了得到每一个的\texttt{ML\_RDES\_SPARSDES}的值，方案如下
\begin{itemize}
	\item 根据\texttt{ML\_RDES\_SPARSDES}值，重新拟合\textrm{ML\_FF}模型
	\item 确定\textrm{MLFF}模型的测试集精度(\textcolor{red}{步骤见见面})
	\item 确定\textrm{MLFF}模型的耗时情况
\end{itemize}
为了获得有意义的耗时结果，需要考虑多个\textrm{MD}模拟步的力计算部分时间并求其平均值，不过，如果将平均值写入文件也会严重降低执行效率，因此该步骤通常可以剔除。计算控制参数\texttt{ML\_OUTBLOCK}和\texttt{ML\_OUTPUT\_MODE}就是控制该部分的。

计算控制稀疏化文件\textrm{INCAR.sparse}

计算执行文件\textrm{INCAR.run}

计算耗时平均取值文件\textrm{INCAR.timing}

\ch{LiH}结构文件\textrm{POSCAR}

\texttt{ML\_AB}文件中包括前面\ch{LiH}第一原理计算数据:~晶格矢量，原子位置，能量，力和应力。

本部分计算中，有两类不同的\textrm{INCAR}文件
\begin{itemize}
	\item \textrm{INCAR.sparse}:~确定三体描述符，通过减少描述符重新拟合力场。这里设置计算控制参数\texttt{ML\_RCUT1}=13.0，因为这个值就是此前得到的优化值
	\item \textrm{INCAR.run}:~对于给定的\textrm{POSCAR}计算能量、力和应力丈量。
\end{itemize}
再次使用前面第一原理计算得到的\texttt{ML\_AB}文件中的数据。

\textrm{POSCAR}文件给定的是\textrm{MD}计算的初始结构，与稀疏化的过程无关。

执行如下的\textrm{Bash}脚本完成稀疏化。

对比前次计算的精确模型和稀疏化之后模型的描述符数目。描述符数目可以在\texttt{ML\_LOGFILE}文件中通过关键词\textrm{NDESC}检索到，执行\textrm{Bash}命令

在\textrm{grep}输出的最后，是描述符的数目。执行以下\textrm{Bash}脚本，对比稀疏化前后的描述符数目

从输出结果可以看到，两种情况下，径向描述符相同，都是24个，角度部分的描述符，稀疏化后减少了一半，两类原子有关的角度部分描述符都从452降为225。

下面，与前面一样，再分析得到力场的情况。通过\textrm{Bash}脚本和\textrm{py4vasp}进行误差分析：

下表给出三种力场的测试集误差

\begin{table}[h!]
    \centering
    \begin{tabular*}{\textwidth}{@{\extracolsep{\fill}}c@{\extracolsep{\fill}}|@{\extracolsep{\fill}}c}
        \toprule
	标准\textrm{MLFF} &\textrm{MLFF}有关\texttt{ML\_RCUT1} &\textrm{MLFF}稀疏化\\
	\midrule
	\textrm{8.5E-3~eV/\AA} &\textrm{7.9E-3~eV/\AA} &\textrm{8.2E-3~eV/\AA}
        \bottomrule
    \end{tabular*}
    \caption{三种力场的误差分析}
\label{Tab:Test-set_Error-in-Sparcsified}
\end{table}

稀疏化使得测试集误差略有增大，但为每一类原子各减少了一半的加度描述符。

随后\textcolor{red}{的相模拟}会评估\textrm{MLFF}模型的模拟能力。这种\textrm{MLFF}力场分析对于计算节点上同时执行的任务数会比较敏感，因为单个节点上任务越多，耗时越长。不过这里讨论的是怎分析\textrm{MLFF}模型的计算性能。
\begin{itemize}
	\item首先，减少计算中的\textrm{I/O}。在\textrm{INCAR.timing}中有如何减少每个\textrm{MD}执行过程中的写磁盘和输出的控制参数方案。为执行计算，将\textrm{INCAR.timing}另存为\textrm{INCAR}，开启\textrm{VASP}计算。
	\item 针对不同的\textrm{MLFF}，执行\textrm{MD}计算，并保存到\textrm{OUTCAR}文件中。
	\item 最后，执行``\textrm{grep LOOP+}''命令，对每个\textrm{MD}模拟步的耗时求\textrm{CPU}耗时的平均值。%如下的输出文件\textrm{OUTCAR.sparse_time}和\textrm{OUTCAR.no-sparse_timing}文件是预先准备好的。
\end{itemize}
注意，这里使用的是较小的第一原理数据(\texttt{ML\_AB}文件)。使用更大的数据库，稀疏化更有效，可以节约更多的时间。系数化最大可达的加速因子由计算控制参数\texttt{ML\_RDES\_SPARSDES}确定。

\subsection{采用优化的\textrm{MLFF}模型完成反应模拟}
根据得到的稀疏化的\textrm{MLFF}模型，来完成反应模拟。因为模型已经经过正确的训练，参数保存在\textrm{ML\_FF}文件中，因此无需对\textrm{INCAR}文件作任何改动。在\textrm{INCAR}文件中，与\textrm{ML\_FF}模型有关的仅有的计算控制参数为\texttt{ML\_LMLFF}=\textrm{TRUE}和\texttt{ML\_MODE}=\textrm{run}。在\textrm{ML\_FF}文件的第一行，记录了构建\textrm{MLFF}模型使用的超参数。而\textrm{ML\_FF}文件的剩下部分采用非格式输出(二进制格式)。采用如下\textrm{Bash}脚本检验\textrm{ML\_FF}模型使用的超参数

输出为

不难看出，优化后的\texttt{ML\_RCUT1}取为13.0，而对于元素\ch{Li}和\ch{H}的描述符数目分别取为$[249,249]$。

使用该\textrm{MLFF}模型完成计算，执行\textrm{MD}模拟计算对分布函数\textrm{(pair-distribution function)}。对分布函数与静态结构因子满足\textrm{Fourier}变换关系。静态结构因子是可以通过散射实验，如\textrm{X-}射线衍射或中子散射实验直接得到。因此通过分布函数，就能通过实验数据来验证\textrm{MLFF}模型的可靠性。%有关结构因此的信息参阅文献\cite{PR95-249_1954}。

这里模拟\ch{LiH}的对分布函数的条件是温度\textrm{450K}，压力\textrm{1bar}，采用\textrm{Langevin}等温条件维持温度和压力。

计算控制文件\textrm{INCAR}

结构文件\textrm{POSCAR}

力场文件\textrm{ML\_FF}

这里的\textrm{POSCAR}文件已经预先弛豫和达到平衡，并包含了起始初速度。

模拟使用之前建立的力场。%在\textrm{Bash}命令下利用\textrm{link}命令，减少存储开销。

\textrm{POTCAR}文件在\textrm{VASP}计算中必须要存在，尽管纯的\textrm{MLFF}模型模拟完用不到\textrm{POTCAR}。所以\textrm{POTCAR}的内容反而无足轻重。

使用以下\textrm{Bash}脚本，执行\textrm{MLFF}模型模拟。

首先作出能量随温度变化的图，确保在\textrm{MD}模拟执行过程中，能量守恒且温度在期望的范围内涨落。

如果能量和温度曲线在一个确定的平均值附近波动，可以认为这个\textrm{MD}模拟得以正确执行。

利用\textrm{Py4vasp}模块作出对分布函数。

从\textrm{OSZICAR}提取能量和温度数据，并作图，采用\textrm{gnuplot}作图命令依次为

如果没有\textrm{Py4vasp}，使命的\textrm{Bash}脚本为\textrm{extract\_PCDAT.sh}，该脚本产生文件\textrm{PCDAT.xy}，因此可以用\textrm{gnuplot}命令绘图。
